\chapter{Outlier detection and treatment}
\label{ch:outliers}

\begin{quote}
\emph{``An outlier is not noise until you have proven it is noise.''}
\end{quote}

% ============================================================
\section{What is an outlier?}

An outlier is an observation that deviates substantially from the pattern established by the majority of the data. Outliers can be:

\begin{itemize}
  \item \textbf{Data entry errors}: a patient weight of 800\,kg (meant 80).
  \item \textbf{Measurement artifacts}: a sensor returning $-999$ during calibration.
  \item \textbf{Genuine extreme values}: a house selling for \$10 million in a suburb where the median is \$500k.
\end{itemize}

\begin{warningbox}
Never remove outliers automatically. Each outlier must be investigated. Removing genuine extreme values because they ``look wrong'' corrupts your analysis. Removing data errors because they ``look unusual'' is correct and necessary.
\end{warningbox}

% ============================================================
\section{Visual detection}

\begin{minted}{python}
import pandas as pd
import matplotlib.pyplot as plt
import seaborn as sns

# FIFA 21 dataset has many messy, extreme values
df = pd.read_csv("fifa21_raw.csv")

# Box plots reveal outliers immediately
fig, axes = plt.subplots(1, 3, figsize=(14, 5))

sns.boxplot(y=df["Age"], ax=axes[0])
axes[0].set_title("Age distribution")

sns.boxplot(y=df["Overall"], ax=axes[1])
axes[1].set_title("Overall rating")

sns.boxplot(y=df["Value(in Euro)"], ax=axes[2])
axes[2].set_title("Market value (Euro)")

plt.tight_layout()
plt.savefig("outlier_boxplots.png", dpi=150)
plt.show()
\end{minted}

\begin{minted}{python}
# Scatter plot: outliers in 2D space
fig, ax = plt.subplots(figsize=(8, 6))
ax.scatter(df["Age"], df["Overall"], alpha=0.3, s=10)
ax.set_xlabel("Age")
ax.set_ylabel("Overall Rating")
ax.set_title("Age vs Overall — look for isolated points")
plt.tight_layout()
plt.show()
\end{minted}

% ============================================================
\section{Statistical methods}

\subsection{Z-score method}

\begin{minted}{python}
import numpy as np
from scipy import stats

# Z-score: number of standard deviations from the mean
z_scores = np.abs(stats.zscore(df["Overall"].dropna()))
threshold = 3
outliers_z = (z_scores > threshold)
print(f"Outliers by z-score (|z| > {threshold}): {outliers_z.sum()}")
\end{minted}

\begin{minted}{python}
# Apply z-score to multiple columns
def detect_zscore_outliers(df, columns, threshold=3):
    """Return a boolean mask of rows with any z-score outlier."""
    mask = pd.Series(False, index=df.index)
    for col in columns:
        z = np.abs(stats.zscore(df[col].dropna()))
        col_mask = pd.Series(False, index=df.index)
        col_mask.iloc[:len(z)] = z > threshold
        mask = mask | col_mask
    return mask

num_cols = ["Age", "Overall", "Potential"]
outlier_mask = detect_zscore_outliers(df, num_cols)
print(f"Rows with at least one outlier: {outlier_mask.sum()}")
\end{minted}

\begin{datatip}
The z-score method assumes a roughly normal distribution. For skewed data (income, house prices, player market values), use the IQR method or log-transform first.
\end{datatip}

\subsection{IQR method}

\begin{minted}{python}
def iqr_bounds(series):
    """Return lower and upper IQR bounds."""
    Q1 = series.quantile(0.25)
    Q3 = series.quantile(0.75)
    IQR = Q3 - Q1
    lower = Q1 - 1.5 * IQR
    upper = Q3 + 1.5 * IQR
    return lower, upper

lower, upper = iqr_bounds(df["Overall"])
outliers_iqr = df[(df["Overall"] < lower) | (df["Overall"] > upper)]
print(f"IQR bounds: [{lower:.1f}, {upper:.1f}]")
print(f"Outliers: {len(outliers_iqr)}")
\end{minted}

\begin{minted}{python}
# IQR report for all numerical columns
def iqr_outlier_report(df, columns):
    """Generate an IQR outlier report."""
    rows = []
    for col in columns:
        lo, hi = iqr_bounds(df[col].dropna())
        n_out = ((df[col] < lo) | (df[col] > hi)).sum()
        rows.append({"column": col, "lower": lo, "upper": hi,
                      "n_outliers": n_out,
                      "pct": round(n_out / len(df) * 100, 1)})
    return pd.DataFrame(rows)

report = iqr_outlier_report(df, ["Age", "Overall", "Potential"])
print(report)
\end{minted}

% ============================================================
\section{Multivariate outlier detection}

\subsection{Mahalanobis distance}

\begin{minted}{python}
from scipy.spatial.distance import mahalanobis
from numpy.linalg import inv

def mahalanobis_outliers(df, columns, threshold=3):
    """Detect multivariate outliers using Mahalanobis distance."""
    data = df[columns].dropna()
    mean = data.mean().values
    cov_matrix = data.cov().values
    cov_inv = inv(cov_matrix)

    distances = data.apply(
        lambda row: mahalanobis(row.values, mean, cov_inv), axis=1
    )
    return distances, distances > threshold

cols = ["Age", "Overall", "Potential"]
distances, mask = mahalanobis_outliers(df, cols)
print(f"Multivariate outliers: {mask.sum()}")
\end{minted}

\begin{preprocesstip}
Mahalanobis distance detects points that are not outliers in any single dimension but are unusual in the \emph{combination} of dimensions. A 40-year-old player is normal. A 95-rated player is normal. A 40-year-old with a 95 rating might be an outlier in the joint space.
\end{preprocesstip}

\subsection{Isolation Forest}

\begin{minted}{python}
from sklearn.ensemble import IsolationForest

# Isolation Forest: tree-based anomaly detection
iso = IsolationForest(contamination=0.05, random_state=42)
num_data = df[["Age", "Overall", "Potential"]].dropna()
labels = iso.fit_predict(num_data)

# -1 = outlier, 1 = inlier
n_outliers = (labels == -1).sum()
print(f"Isolation Forest outliers: {n_outliers} "
      f"({n_outliers / len(num_data) * 100:.1f}%)")
\end{minted}

% ============================================================
\section{Domain-specific rules}

\begin{minted}{python}
# Air Quality dataset: domain knowledge matters
df_air = pd.read_csv("AirQualityUCI.csv", sep=";", decimal=",")

# The dataset uses -200 as a sentinel for missing sensor data
print(f"Values == -200: {(df_air == -200).sum().sum()}")

# Replace sentinel with NaN
df_air = df_air.replace(-200, np.nan)

# Domain rule: CO concentration > 20 mg/m3 is physically implausible
# at urban monitoring stations
co_col = "CO(GT)"
implausible = df_air[co_col] > 20
print(f"Implausible CO readings: {implausible.sum()}")
df_air.loc[implausible, co_col] = np.nan
\end{minted}

\begin{datatip}
Domain rules are the most reliable outlier detection method. A z-score does not know that a human cannot weigh 800\,kg. Always consult domain experts or reference ranges when available.
\end{datatip}

% ============================================================
\section{Treatment strategies}

\begin{minted}{python}
# Strategy 1: Remove outlier rows
df_clean = df[~outlier_mask].copy()

# Strategy 2: Cap (winsorize) at bounds
from scipy.stats import mstats
df["Overall_winsorized"] = mstats.winsorize(df["Overall"], limits=[0.01, 0.01])

# Strategy 3: Replace with NaN and impute later
df.loc[outlier_mask, "Overall"] = np.nan

# Strategy 4: Log transform to reduce skew
df["Value_log"] = np.log1p(df["Value(in Euro)"].clip(lower=0))

# Strategy 5: Robust scaling (less sensitive to outliers)
from sklearn.preprocessing import RobustScaler
scaler = RobustScaler()  # uses median and IQR instead of mean and std
df[["Overall_robust"]] = scaler.fit_transform(df[["Overall"]])
\end{minted}

\begin{minted}{python}
# Compare distributions before and after winsorization
fig, axes = plt.subplots(1, 2, figsize=(12, 4))

df["Overall"].hist(bins=50, ax=axes[0], alpha=0.7)
axes[0].set_title("Original")

df["Overall_winsorized"].hist(bins=50, ax=axes[1], alpha=0.7, color="green")
axes[1].set_title("Winsorized (1% each tail)")

plt.tight_layout()
plt.show()
\end{minted}

% ============================================================
\section{Exercises}

\begin{exercise}
\begin{enumerate}
  \item Load the FIFA 21 dataset. Create box plots for five numerical columns. Identify which columns have the most outliers using the IQR method.
  \item For the Air Quality UCI dataset, replace all $-200$ sentinel values with NaN. Then apply domain-specific rules to flag implausible values for at least two pollutant columns.
  \item Compare z-score and IQR outlier counts for a skewed column (e.g., player market value). Which method is more appropriate and why?
  \item Use Isolation Forest on the Melbourne Housing dataset (Price, Rooms, Distance, Landsize). Visualize the detected outliers on a scatter plot of Price vs.\ Landsize.
  \item Implement a function \texttt{treat\_outliers(series, method)} that supports three methods: ``remove'', ``cap'', and ``log''. Test it on a column with known outliers.
\end{enumerate}
\end{exercise}

% ============================================================
\section{Chapter summary}

\begin{itemize}
  \item Outliers can be errors, artifacts, or genuine extremes --- investigation is mandatory.
  \item Visual methods (box plots, scatter plots) provide immediate intuition about outlier presence.
  \item Z-score works for normally distributed data; IQR is robust to skewed distributions.
  \item Mahalanobis distance and Isolation Forest detect multivariate outliers invisible to univariate methods.
  \item Domain-specific rules (physical limits, sentinel values) are the most reliable detection method.
  \item Treatment options include removal, capping (winsorization), log transformation, and robust scaling.
\end{itemize}
